\expandafter\gdef\csname SNside1\endcsname{\node[anchor=north west,inner sep=0pt,text width=48mm,font=\SideFont,align=left] (SNsidebox) at ([xshift=138mm,yshift=-402.91pt]current page.north west) {\textcolor{SNink}{\textbf{课程连接 · 链路反馈只看一跳}}\par\vspace{1.5mm}操作系统中的缓冲区保护解决本机生产者与消费者速度不匹配；链路流控也限制发送量以保护接收缓存。但链路ACK只能覆盖当前相邻节点，不能像端到端协议一样替后续整条路径承诺可靠。};\pgfextracty{\SNsideheight}{\pgfpointdiff{\pgfpointanchor{SNsidebox}{south}}{\pgfpointanchor{SNsidebox}{north}}}\typeout{SNHEIGHT:1:0:\the\SNsideheight}}
\expandafter\gdef\csname SNside2\endcsname{\node[anchor=north west,inner sep=0pt,text width=48mm,font=\SideFont,align=left] (SNsidebox) at ([xshift=138mm,yshift=-301.44pt]current page.north west) {\textcolor{SNink}{\textbf{最小反例 · 合法码字之间仍能跳转}}\par\vspace{1.5mm}取生成式11，000与011均能被它整除。合法串000若恰好翻转后两位成为011，检验余数仍为0。这个两位翻转例子直接排除``CRC通过就证明没有任何错误''的绝对断言。\par\smallskip\begin{tikzpicture}[x=1mm,y=1mm,>=Stealth,node distance=5mm]
\node[draw=SNink!50,fill=SNblue,rounded corners=1mm,text width=34mm,align=center,inner sep=1.2mm,font=\fontsize{7.5}{9}\selectfont] (m0) at (0,0) {000合法};
\node[draw=SNink!50,fill=SNblue,rounded corners=1mm,text width=34mm,align=center,inner sep=1.2mm,font=\fontsize{7.5}{9}\selectfont] (m1) at (0,-11) {翻后两位};
\node[draw=SNink!50,fill=SNblue,rounded corners=1mm,text width=34mm,align=center,inner sep=1.2mm,font=\fontsize{7.5}{9}\selectfont] (m2) at (0,-22) {011仍合法};
\draw[->,SNink] (m0) -- (m1);
\draw[->,SNink] (m1) -- (m2);
\end{tikzpicture}
\par{\fontsize{7.5}{9}\selectfont 生成式11下，错误仍可能通过检验}};\pgfextracty{\SNsideheight}{\pgfpointdiff{\pgfpointanchor{SNsidebox}{south}}{\pgfpointanchor{SNsidebox}{north}}}\typeout{SNHEIGHT:2:2:\the\SNsideheight}}
\expandafter\gdef\csname SNside3\endcsname{\node[anchor=north west,inner sep=0pt,text width=48mm,font=\SideFont,align=left] (SNsidebox) at ([xshift=138mm,yshift=-95.00pt]current page.north west) {\textcolor{SNink}{\textbf{不变量 · 填充只删自己插入的位}}\par\vspace{1.5mm}数据中五个连续1之后插入的0在接收端删除，原始数据的0必须留下。遇到数据111110时线上会出现1111100；如果把两个0都删掉，成帧机制反而破坏了上层数据。};\pgfextracty{\SNsideheight}{\pgfpointdiff{\pgfpointanchor{SNsidebox}{south}}{\pgfpointanchor{SNsidebox}{north}}}\typeout{SNHEIGHT:3:1:\the\SNsideheight}}
\expandafter\gdef\csname SNside4\endcsname{\node[anchor=north west,inner sep=0pt,text width=48mm,font=\SideFont,align=left] (SNsidebox) at ([xshift=138mm,yshift=-95.00pt]current page.north west) {\textcolor{SNink}{\textbf{不变量 · 综合征须固定编号}}\par\vspace{1.5mm}位置1、2、4的检验位与书写方向是两套约定。先给每位编号，再按二进制位置选校验集合；单错时综合征才指向位置。若有双错，直接把综合征当单错位置翻转可能产生新的错误。};\pgfextracty{\SNsideheight}{\pgfpointdiff{\pgfpointanchor{SNsidebox}{south}}{\pgfpointanchor{SNsidebox}{north}}}\typeout{SNHEIGHT:4:3:\the\SNsideheight}}
\expandafter\gdef\csname SNside5\endcsname{\node[anchor=north west,inner sep=0pt,text width=48mm,font=\SideFont,align=left] (SNsidebox) at ([xshift=138mm,yshift=-95.00pt]current page.north west) {\textcolor{SNink}{\textbf{最小反例 · 没收到ACK也可能交付过}}\par\vspace{1.5mm}数据0已被接收并交付，返回ACK0却丢失。发送方超时重发0，接收方必须重复确认而不重复交付。若没有序号，接收端就无法区分这次重传与内容恰好相同的新数据。};\pgfextracty{\SNsideheight}{\pgfpointdiff{\pgfpointanchor{SNsidebox}{south}}{\pgfpointanchor{SNsidebox}{north}}}\typeout{SNHEIGHT:5:4:\the\SNsideheight}}
\expandafter\gdef\csname SNside6\endcsname{\node[anchor=north west,inner sep=0pt,text width=48mm,font=\SideFont,align=left] (SNsidebox) at ([xshift=138mm,yshift=-95.00pt]current page.north west) {\textcolor{SNink}{\textbf{错项排除 · 累计确认覆盖缺失的ACK}}\par\vspace{1.5mm}GBN收到ACK3通常已证明本轮0至3全部按序接收，即使ACK1没有单独到达。重传集合由累计左沿决定，不能按``哪些ACK没看见''逐个列出；换成SR独立确认时结论又会改变。};\pgfextracty{\SNsideheight}{\pgfpointdiff{\pgfpointanchor{SNsidebox}{south}}{\pgfpointanchor{SNsidebox}{north}}}\typeout{SNHEIGHT:6:6:\the\SNsideheight}}
\expandafter\gdef\csname SNside7\endcsname{\node[anchor=north west,inner sep=0pt,text width=48mm,font=\SideFont,align=left] (SNsidebox) at ([xshift=138mm,yshift=-95.00pt]current page.north west) {\textcolor{SNink}{\textbf{最小反例 · 缓存不等于越过缺口}}\par\vspace{1.5mm}接收窗口允许2和3，3先到只能缓存，左沿仍停在2；2到达后才一起交付2、3。画出这两个时刻就能看出``收到高序号帧后立刻把窗口移到它后面''为什么会丢掉尚未收到的数据。};\pgfextracty{\SNsideheight}{\pgfpointdiff{\pgfpointanchor{SNsidebox}{south}}{\pgfpointanchor{SNsidebox}{north}}}\typeout{SNHEIGHT:7:7:\the\SNsideheight}}
\expandafter\gdef\csname SNside8\endcsname{\node[anchor=north west,inner sep=0pt,text width=48mm,font=\SideFont,align=left] (SNsidebox) at ([xshift=138mm,yshift=-95.00pt]current page.north west) {\textcolor{SNink}{\textbf{课程连接 · 正交像向量基的分解}}\par\vspace{1.5mm}理想CDMA中，与目标码正交的其他用户贡献在内积中为0；目标码与自身内积为码长。恢复的是目标发送者的数据，接收站名称不会改变要选哪条码；同步与功率假设也是消去干扰的条件。};\pgfextracty{\SNsideheight}{\pgfpointdiff{\pgfpointanchor{SNsidebox}{south}}{\pgfpointanchor{SNsidebox}{north}}}\typeout{SNHEIGHT:8:9:\the\SNsideheight}}
\expandafter\gdef\csname SNside9\endcsname{\node[anchor=north west,inner sep=0pt,text width=48mm,font=\SideFont,align=left] (SNsidebox) at ([xshift=138mm,yshift=-95.00pt]current page.north west) {\textcolor{SNink}{\textbf{不变量 · 最短帧要覆盖往返}}\par\vspace{1.5mm}最远站点开始发送的时刻可以恰在本端信号到达前，冲突反馈还要传播回来。最短发送时间约束的是两倍单向传播时间；把最短帧缩短800 bit时，距离变化必须再除以往返系数2。};\pgfextracty{\SNsideheight}{\pgfpointdiff{\pgfpointanchor{SNsidebox}{south}}{\pgfpointanchor{SNsidebox}{north}}}\typeout{SNHEIGHT:9:10:\the\SNsideheight}}
\expandafter\gdef\csname SNside10\endcsname{\node[anchor=north west,inner sep=0pt,text width=48mm,font=\SideFont,align=left] (SNsidebox) at ([xshift=138mm,yshift=-95.00pt]current page.north west) {\textcolor{SNink}{\textbf{错项排除 · NAV只预约未来}}\par\vspace{1.5mm}收到CTS的一刻，RTS、CTS和前面的等待都已过去。此后还剩的SIFS、DATA、SIFS、ACK才构成基本预约余量。不同题目从不同事件起计时，应先画起点，不能只凭协议名背加法项。};\pgfextracty{\SNsideheight}{\pgfpointdiff{\pgfpointanchor{SNsidebox}{south}}{\pgfpointanchor{SNsidebox}{north}}}\typeout{SNHEIGHT:10:11:\the\SNsideheight}}
\expandafter\gdef\csname SNside11\endcsname{\node[anchor=north west,inner sep=0pt,text width=48mm,font=\SideFont,align=left] (SNsidebox) at ([xshift=138mm,yshift=-95.00pt]current page.north west) {\textcolor{SNink}{\textbf{不变量 · 帧与线上周期不同}}\par\vspace{1.5mm}无标签最短MAC帧64 B，若再计8 B前导部分与96 bit时间的帧间隔，一个发送周期相当于84 B。间隔是时间而非额外字段；IP总长度、MAC帧长与线上周期必须各自保持边界。};\pgfextracty{\SNsideheight}{\pgfpointdiff{\pgfpointanchor{SNsidebox}{south}}{\pgfpointanchor{SNsidebox}{north}}}\typeout{SNHEIGHT:11:12:\the\SNsideheight}}
\expandafter\gdef\csname SNside12\endcsname{\node[anchor=north west,inner sep=0pt,text width=48mm,font=\SideFont,align=left] (SNsidebox) at ([xshift=138mm,yshift=-95.00pt]current page.north west) {\textcolor{SNink}{\textbf{错项排除 · 当前跳发送者未必是源}}\par\vspace{1.5mm}AP向接收站转发来自远端主机的数据时，无线帧地址2标记当前发送AP，地址3保留原始源。把三个地址都按以太网``源和目的''两字段硬套，会把无线跳的接口和端到端主机混在一起。};\pgfextracty{\SNsideheight}{\pgfpointdiff{\pgfpointanchor{SNsidebox}{south}}{\pgfpointanchor{SNsidebox}{north}}}\typeout{SNHEIGHT:12:13:\the\SNsideheight}}
\expandafter\gdef\csname SNside13\endcsname{\node[anchor=north west,inner sep=0pt,text width=48mm,font=\SideFont,align=left] (SNsidebox) at ([xshift=138mm,yshift=-95.00pt]current page.north west) {\textcolor{SNink}{\textbf{最小反例 · 端口并行不隔离广播}}\par\vspace{1.5mm}A\ensuremath{\to}B与C\ensuremath{\to}D能在四端口交换机上并行，只说明出口互不冲突。一帧广播仍可传播至其余允许端口，所以普通二层交换的独立端口不能自动构成多个广播域；VLAN边界需要另行配置。};\pgfextracty{\SNsideheight}{\pgfpointdiff{\pgfpointanchor{SNsidebox}{south}}{\pgfpointanchor{SNsidebox}{north}}}\typeout{SNHEIGHT:13:14:\the\SNsideheight}}
\expandafter\gdef\csname SNside14\endcsname{\node[anchor=north west,inner sep=0pt,text width=48mm,font=\SideFont,align=left] (SNsidebox) at ([xshift=138mm,yshift=-95.00pt]current page.north west) {\textcolor{SNink}{\textbf{不变量 · 学源查目的}}\par\vspace{1.5mm}帧从端口1带着源A进入，才有证据把A记在端口1；目的B未知时只能泛洪寻找。用目的地址学习会把尚未找到的B误记到发送者位置，下一次转发将沿错误出口走。};\pgfextracty{\SNsideheight}{\pgfpointdiff{\pgfpointanchor{SNsidebox}{south}}{\pgfpointanchor{SNsidebox}{north}}}\typeout{SNHEIGHT:14:16:\the\SNsideheight}}
\expandafter\gdef\csname SNside15\endcsname{\node[anchor=north west,inner sep=0pt,text width=48mm,font=\SideFont,align=left] (SNsidebox) at ([xshift=138mm,yshift=-95.00pt]current page.north west) {\textcolor{SNink}{\textbf{辨析 · 确认与发送分开追踪}}\par\vspace{1.5mm}捎带确认号说明期望收到的下一帧；发送序号记录本方发送。先画两个方向的序号轴，再处理丢失、超时和模8回绕。};\pgfextracty{\SNsideheight}{\pgfpointdiff{\pgfpointanchor{SNsidebox}{south}}{\pgfpointanchor{SNsidebox}{north}}}\typeout{SNHEIGHT:15:21:\the\SNsideheight}}
\expandafter\gdef\csname SNside16\endcsname{\node[anchor=north west,inner sep=0pt,text width=48mm,font=\SideFont,align=left] (SNsidebox) at ([xshift=138mm,yshift=-95.00pt]current page.north west) {\textcolor{SNink}{\textbf{条件 · 先固定书写方向}}\par\vspace{1.5mm}位号1、2、4是校验位置；从H7向H1书写与从H1向H7书写不能混用。};\pgfextracty{\SNsideheight}{\pgfpointdiff{\pgfpointanchor{SNsidebox}{south}}{\pgfpointanchor{SNsidebox}{north}}}\typeout{SNHEIGHT:16:17:\the\SNsideheight}}
\expandafter\gdef\csname SNside17\endcsname{\node[anchor=north west,inner sep=0pt,text width=48mm,font=\SideFont,align=left] (SNsidebox) at ([xshift=138mm,yshift=-95.00pt]current page.north west) {\textcolor{SNink}{\textbf{方法 · 发送时间也占周期}}\par\vspace{1.5mm}首个确认返回周期含本帧发送、往返传播和ACK发送。};\pgfextracty{\SNsideheight}{\pgfpointdiff{\pgfpointanchor{SNsidebox}{south}}{\pgfpointanchor{SNsidebox}{north}}}\typeout{SNHEIGHT:17:18:\the\SNsideheight}}
\expandafter\gdef\csname SNside18\endcsname{\node[anchor=north west,inner sep=0pt,text width=48mm,font=\SideFont,align=left] (SNsidebox) at ([xshift=138mm,yshift=-95.00pt]current page.north west) {\textcolor{SNink}{\textbf{课程连接 · 先LCP后NCP}}\par\vspace{1.5mm}先把链路配置好，再配置承载的网络层协议，之后传IP分组。};\pgfextracty{\SNsideheight}{\pgfpointdiff{\pgfpointanchor{SNsidebox}{south}}{\pgfpointanchor{SNsidebox}{north}}}\typeout{SNHEIGHT:18:20:\the\SNsideheight}}
\expandafter\gdef\csname SNside19\endcsname{\node[anchor=north west,inner sep=0pt,text width=48mm,font=\SideFont,align=left] (SNsidebox) at ([xshift=138mm,yshift=-95.00pt]current page.north west) {\textcolor{SNink}{\textbf{课程连接 · 先LCP后NCP}}\par\vspace{1.5mm}先把链路配置好，再配置承载的网络层协议，之后传IP分组。};\pgfextracty{\SNsideheight}{\pgfpointdiff{\pgfpointanchor{SNsidebox}{south}}{\pgfpointanchor{SNsidebox}{north}}}\typeout{SNHEIGHT:19:20:\the\SNsideheight}}
\AddToShipoutPictureFG{\begin{tikzpicture}[remember picture,overlay]\ifcsname SNside\thepage\endcsname\csname SNside\thepage\endcsname\fi\end{tikzpicture}}